\section{结论}

本文提出 KCH-MedRank：一种知识约束学习排序框架，用于生物医学证据检索与重排序。该方法保持第一阶段检索的实用性，将检索和生物医学语义信号与 MeSH、实体和关系特征结合，并可选加入局部图或超图结构特征以增强可解释性。扁平医学知识 LTR 基线和成对图分解作为受控结构基线。比较表明，扁平生物医学知识特征解释了大部分总体收益；超图在早期排序位置（$k{=}5$）提供适度收益，但在 $k{=}10$ 上未显著优于成对图或扁平医学知识基线。

BioASQ 结果表明，该方法相对于原始与增强 Hybrid RRF 均取得了更强的 top-10 证据检索效果，并且相对于增强 Hybrid RRF 的多个 top-10 检索指标达到统计显著。在同一增强候选池上，相对于真实 MedCPT Cross-Encoder 基线，最有力的结果是 Recall@10 的显著提升；MRR@10 与 nDCG@10 为正向但未显著的趋势。受控重排序阶段计时进一步显示，KCH-MedRank 使用预计算的稠密生物医学语义特征，而不是 Cross-Encoder 特征；在不计候选生成、一次性模型加载和离线训练时，KCH-MedRank 在留出测试候选池上比 MedCPT Cross-Encoder 快 18.6 倍。

结构消融为核心医学知识贡献提供了证据。不含医学知识标注的局部超图引入结构噪声，低于扁平医学知识 LTR 基线。引入生物医学实体特征、精确与层级感知 MeSH 特征及辅助关系特征后，相对于无医学知识超图变体逆转了这一效应。这证明医学知识约束对于局部结构有用而非噪声是必要的。同时，完整模型在主要 top-10 深度上并未显著优于扁平医学知识 LTR 基线，因此可支持的贡献是一个具有互补结构特征的白盒知识约束证据选择框架，而不是独立的超图突破。外部 BEIR NFCorpus retrieval-only 诊断支持监督检索特征重排序的稳健性，同时也表明不能在未进行数据集特定构建的情况下声称完整医学知识特征栈可直接迁移。

附录中的 PubMedQA Qwen3-8B pilot 表明，在固定 prompt 和解码设置下，证据排序会影响下游答案选择，但它仍是有限诊断，而不是主要生成主张。总体而言，结果支持将 KCH-MedRank 作为循证医学问答中的白盒证据选择模块。该方法不应被视为经过临床验证的答案生成或决策支持系统。
